\section{The Tetra codebook}
\label{sec:tetra}

\subsection{The lattice and the quantizer}

We use the integer form of the Leech lattice. A point $y\in\mathbb{Z}^{24}$
lies in $\Lambda_{24}$, up to the global scale the quantizer fixes, when there
are a bit $p$, a codeword $c$ of the binary Golay code $\mathcal{G}_{24}$ and
integers $k_j$ with
\begin{equation}
y_j = p + 2c_j + 4k_j,\qquad \textstyle\sum_j k_j \equiv p \pmod 2 .
\label{eq:leech}
\end{equation}
Two constraints cut $\mathbb{Z}^{24}$ down to the lattice
(Figure~\ref{fig:geom}a): $c$ must be one of the $2^{12}$ Golay codewords among
the $2^{24}$ bit patterns, and one parity couples all 24 coordinates. A decoder
has to respect both.

The encoder is the shape-gain scheme of \citet{llvq2026} inside a GPTQ-style
loop~\citep{gptq2023}, after a rotation of the input~\citep{quip2023,quarot2024}.
For each block of 24 rotated weights it finds the nearest point of a norm-capped
region of $\Lambda_{24}$ and reconstructs
\begin{equation}
\hat{w}_j = \frac{y_j}{\lVert y\rVert}\,\mathrm{gain}[g]\,\sigma_{\mathrm{row}},
\qquad \lVert y\rVert=\sqrt{16m},
\label{eq:recon}
\end{equation}
with $m$ the shell index, $g$ a gain bit and $\sigma_{\mathrm{row}}$ one f32
scale per output row. We change two things in that pipeline: the region the
encoder searches (\S\ref{sec:kept}) and the label it writes. The loop, the
rotation, the gain values and the row scales stay as they were, and the
encoder still runs once, offline, on a CPU.

\subsection{Three octads and a trellis}

An octad is a Golay codeword of weight 8, and three disjoint octads cover the
24 coordinates. We fix one such triple and read the coordinates in that order,
octad by octad. In this order the code splits into sections: cutting after
coordinates 8 and 16 leaves a trellis with 64 states at each cut
(Figure~\ref{fig:geom}b), the fewest any three-way cut of $\mathcal{G}_{24}$
reaches. The trellis has $64\times2\times16\times2 = 4096 = |\mathcal{G}_{24}|$
paths, so every codeword is one path and every path is a codeword. The count is
checked when the tables are built.

\begin{figure}[t]
\centering
% ---- (a) the two constraints of Equation (1) ----------------------------
\begin{tikzpicture}[x=0.42cm, y=0.52cm]
\node[font=\scriptsize, anchor=east] at (-0.4,3.4) {$y_j\;=$};
\node[font=\scriptsize, anchor=east] at (-0.4,2.4) {$p\;+$};
\node[font=\scriptsize, anchor=east] at (-0.4,1.4) {$2c_j\;+$};
\node[font=\scriptsize, anchor=east] at (-0.4,0.4) {$4k_j$};
\foreach \i in {0,...,23}{
  \draw[okGray, fill=black!7] (\i,3.15) rectangle ++(0.8,0.5);
  \draw[okGray] (\i,2.15) rectangle ++(0.8,0.5);
  \draw[okGray] (\i,1.15) rectangle ++(0.8,0.5);
  \draw[okGray] (\i,0.15) rectangle ++(0.8,0.5);}
\foreach \i in {0,2,3,4,7,10,12,13,15,17,18,20,23}{
  \draw[okGray, fill=okBlue!35] (\i,1.15) rectangle ++(0.8,0.5);}
\node[font=\scriptsize, anchor=north west] at (0,-0.25)
  {(i)\; $c\in\mathcal{G}_{24}$:\; $2^{12}$ of the $2^{24}$ bit patterns};
\node[font=\scriptsize, anchor=north west] at (0,-1.05)
  {(ii)\; $\textstyle\sum_j k_j \equiv p \pmod 2$:\; one parity over all 24 coordinates};
\end{tikzpicture}\\[6pt]
{\footnotesize (a) $\mathbb{Z}^{24}$ under two global constraints}\\[10pt]
% ---- (b) the same code in octad order ------------------------------------
\begin{tikzpicture}[x=1cm, y=1cm]
\node[box, font=\scriptsize] (s0) at (0.5,1.3) {$1$};
\node[mem, font=\scriptsize, minimum width=9mm, minimum height=13mm]
  (s8) at (4.2,1.3) {$64$};
\node[mem, font=\scriptsize, minimum width=9mm, minimum height=13mm]
  (s16) at (8.2,1.3) {$64$};
\node[box, font=\scriptsize] (se) at (11.9,1.3) {$1$};
\draw[flow] (s0) -- node[above, font=\scriptsize]{2 prefixes}
                   node[below, font=\scriptsize]{$b_1$: 1 bit} (s8);
\draw[flow] (s8) -- node[above, font=\scriptsize]{16 branches}
                    node[below, font=\scriptsize]{$b_2$: 4 bits} (s16);
\draw[flow] (s16) -- node[above, font=\scriptsize]{2 suffixes}
                     node[below, font=\scriptsize]{$b_3$: 1 bit} (se);
\node[font=\scriptsize, anchor=north, align=center] at (4.2,0.45)
  {$s_8$: the 6 bits\\the word stores};
\node[font=\scriptsize, anchor=north, align=center] at (8.2,0.45)
  {$s_{16}$: not stored,\\returned with the byte};
\node[font=\scriptsize, anchor=north] at (6.2,-0.55)
  {$64\times2\times16\times2 \;=\; 4096 \;=\; |\mathcal{G}_{24}|$};
\end{tikzpicture}\\[4pt]
{\footnotesize (b) the same code in octad order}
\caption{The object and the reordering. (a) Equation~\eqref{eq:leech}, with the
shaded cells marking $c_j=1$. Three independent 8-dimensional quantizers
would respect neither constraint. (b) In octad order the same code is a
trellis, and a codeword is a path through it.}
\label{fig:geom}
\end{figure}

\subsection{State, rank vectors, and one table for three sections}
\label{sec:state}

The word carries eight bits of state: the trellis state $s_8$ (6 bits), the
shared parity $p$, and $r$, the parity of $\sum k_j$ over section 1. What $s_8$
holds is exactly what section 2 needs to know about section 1. So
$\Lambda_{24}$ adds two bits to the code's own state and no more:
$64\times4 = 256$ lattice states at a cut.

The $k$-parity of each section is never stored as a field. Section 1's is $r$.
Section 2's is $\delta$, read off which half of the shared table its index
falls in. Section 3's is
\begin{equation}
r_3 = p \oplus r \oplus \delta ,
\label{eq:xor}
\end{equation}
so the parity constraint of Equation~\eqref{eq:leech} costs one XOR and no bits.

Inside a section, the residues mod 4 are fixed once $p$ and the pattern byte
are known, $o_j = p + 2c_j$, so coordinate $j$ can only take values in
$o_j + 4\mathbb{Z}$. We store its \emph{rank} in that progression rather than
its value, counting outward from zero: for $o=0$ the values are
$0, +4, -4, +8, \dots$, for $o=2$ they are $+2, -2, +6, \dots$, for $o=1$
$+1, -3, +5, \dots$, and for $o=3$ $-1, +3, -5, \dots$. A section is then a rank
vector $\rho \in \{0,\dots,7\}^8$: four bits per coordinate, one 32-bit word.

Two properties let a single table serve every section of every block. First,
with $\mathrm{cost}(\rho)=\sum_j(2\rho_j+1)^2$, the order of rank vectors by
cost does not depend on $p$ or on the pattern byte, so ``keep the 2{,}048
cheapest'' is one list rather than $4^8$ of them. Second,
$\mathrm{cls}(\rho)=\sum_j[\rho_j\in\{1,2\}] \bmod 2$ equals the parity of
$\sum_j k_j$ for every pattern, so a section's parity class is an address
rather than a field. The table holds 4{,}096 words: the 2{,}048 cheapest rank
vectors of class 0, then the 2{,}048 cheapest of class 1, with eleven bits
addressing a half. The middle section reads the same rows in a third, mixed
order, the 2{,}048 cheapest overall, which is $N_0 = 1240$ rows of class 0
followed by 808 of class 1. Its index therefore gives the $\delta$ that
Equation~\eqref{eq:xor} needs.

\subsection{The word and the decode}
\label{sec:decode}

Figure~\ref{fig:word} shows the resulting 48-bit word next to the unfolded
record it replaces. Every field range is a power of two, so any 48-bit value is
a legal word and decodes to a point of $\Lambda_{24}$. A benchmark can therefore
stream random words and check every decode against a reference.
Figure~\ref{fig:decode} is the decode: six loads per block, three of them on a
dependent chain.

\begin{figure}[t]
\centering
\includegraphics{fig_word.pdf}
\caption{The record of our earlier layout and the \Tetra{} word, at one bit
scale, with the word magnified below so its ten fields can be named.
\Planes{} is what the 47-bit index becomes after unfolding: a 9-bit class, a
gain bit, a sign mask and three bit planes, 112 bits at a 14-byte stride.
\Tetra{} is what the file stores and what the kernel reads: 48 bits at a
6-byte stride, decoded by six loads from 18{,}816~B of tables, 16~KiB of which
hold the rank rows. Purple is trellis state, grey the edge choices, blue the
row indices into the shared table, red the gain bit; hatching is padding. Both
carry 2.000 bits of code per weight.}
\label{fig:word}
\end{figure}

\begin{figure}[t]
\centering
\begin{tikzpicture}[x=1cm, y=1cm]
\node[font=\footnotesize\bfseries, anchor=south] at (4.8,2.62)
  {one block: six loads, three of them dependent};
\node[box, font=\scriptsize, text width=2.7cm] (w) at (4.8,2.1)
  {the 48-bit word};
\node[mem, font=\scriptsize, text width=4.2cm, align=left] (L) at (2.4,0.95)
  {\textbf{trellis chain}, 2\,304~B\\[1pt]
   \texttt{prefixes}$[2s_8{+}b_1]$\\
   \texttt{branches}$[16s_8{+}b_2]$\\
   \texttt{suffixes}$[2s_{16}{+}b_3]$\\[1pt]
   \emph{chained: each read needs\\the previous one}};
\node[mem, font=\scriptsize, text width=4.2cm, align=left] (R) at (7.2,0.95)
  {\textbf{rank table}, 16~KiB\\[1pt]
   \texttt{rows}$[2048r{+}i_1]$\\
   \texttt{rows}$[\mathrm{mix}(i_2)]$\\
   \texttt{rows}$[2048r_3{+}i_3]$\\[1pt]
   \emph{independent: issued together}};
\node[box, font=\scriptsize, text width=7.2cm] (y) at (4.8,-0.85)
  {24 values $y_j$, then $\times\,\mathrm{gain}[g] \times 1/\!\sqrt{16m}$};
\draw[flow] (w.south west) -- (L.north);
\draw[flow] (w.south east) -- (R.north);
\draw[flow] (L.south) -- (y.north west);
\draw[flow] (R.south) -- (y.north east);
\end{tikzpicture}
\caption{The decode. The three byte lookups recover the Golay pattern bytes and
must run in order, since each index depends on the state the previous read
returned. The three row reads depend only on the word and are issued together.
The shell index $m$ is recomputed from the decoded values, not stored, and
$r_3 = p \oplus r \oplus \delta$ by Equation~\eqref{eq:xor}.}
\label{fig:decode}
\end{figure}

The tables total 18{,}816~B: the rank table, 4{,}096 words or 16{,}384~B; the
branch table, 1{,}024 16-bit entries or 2{,}048~B; two byte tables of 128~B;
and 32 floats holding $1/\sqrt{16m}$, with entry 0 set to 0 so that the origin,
a legal word, decodes to zeros without a branch. The bound $m\le27$ comes from
sweeping all 4{,}096 rows against both parities and both residues, which gives
$\max|y_j|=10$.

\subsection{What the construction costs}
\label{sec:kept}

\Tetra{} is not a repacking of the earlier codebook. It is a different codebook
on the same lattice. Keeping the 2{,}048 cheapest rank vectors bounds each
section on its own, so the encoder searches a product of three 8-dimensional
regions where \Planes{} searched a 24-dimensional ball. That is what makes the
six-load decode possible, and it is what the format pays.

We estimate the price rather than measure it. The best shaping gain of a
product of three 8-dimensional balls is 0.7292~dB, against 1.0958~dB for the
24-dimensional ball. The difference, 0.3666~dB, is 8.81\,\% of mean squared
error at our operating point. Our regions are neither balls nor continuous,
since they are rank-truncated and constrained by the trellis, so these figures
are a geometric reference and not the cost of this codebook. One related effect
is measured: ranking rank vectors by $\mathrm{cost}(\rho)$ rather than by exact
squared distance loses 0.6 points of retention against an exact enumeration
per section.

What the construction keeps is the lattice. The decoded point satisfies both
constraints of Figure~\ref{fig:geom}: Golay membership, because the pattern
bytes come from a trellis path, and the $\sum k$ parity, by
Equation~\eqref{eq:xor}. Three independent $E_8$ quantizers would keep neither,
and the lattice alone is worth $G(E_8)/G(\Lambda_{24}) = 0.071682/0.065771 =
1.0899$, that is 9.0\,\% of mean squared error~\citep{conway1999sphere}.

On the same model and the same harness, bare \Tetra{} scores 2.10 MMLU points
below \Planes{} on the 2{,}280-question sample and has a 4.64\,\% lower
perplexity. Re-encoding one model with three different calibration draws moves
MMLU by 2.92 points, more than the difference between the two codebooks, so
this comparison does not separate the codebook from the calibration draw.
